\begin{table*}[t]
\centering
\small
\begin{tabular}{lp{4cm}p{6cm}}
\toprule
Exp. & Title & Headline result \\
\midrule
E7 & Temporal security protocol -- verification suite & All hash/key-lifecycle/AAD tests passed (Tests A-F) \\
E9 & Retrospective breach attack (5 conditions x 1000 trials) & 0/5000 successful plaintext recoveries (0.00\%) \\
E10 & Temporal window (Tr) sweep -- simulated client latency & Tr=120s -> 100\% round completion; Tr>=300s recommended for operational margin \\
E14 & Security-layer scalability (K=3,5,10,20 simulated clients) & \textbf{UNRESOLVED:} two runs on file --- JSON (run A): K=3$\to$179.96\,ms / K=20$\to$752.79\,ms; stdout log (run B): K=3$\to$231.88\,ms / K=20$\to$903.64\,ms. Do not cite without a confirmed re-run. \\
E15 & Fault & late-client robustness (10 scenarios) & 10/10 scenarios passed \\
\bottomrule
\end{tabular}
\caption{Security and robustness experiments. See \texttt{Results\_New/results/MASTER\_RESULTS\_SECURITY.json} for full provenance notes. E9's confidence-interval methodology has been resolved (2026-09-26): report the per-condition bound ($\le 0.30\%$, rule-of-three, $n=1000$ per attack condition), not the pooled $n=5000$ figure. Two items require attention before submission: (1) E10's committed log does not currently evidence the committed result (numbers retained from JSON as authoritative per project decision, but no clean re-run log exists); (2) E14 has an unresolved discrepancy between \texttt{e14\_scalability\_results.json} (run~A) and \texttt{e14\_stdout.log} (run~B) --- both sets of numbers are retained in \texttt{Results\_New/E14/E14\_RESULTS.md} and neither has been declared authoritative.}
\label{tab:master-results-security}
\end{table*}
